\begin{proof}
We prove $\mathrm{pr}\text{-}\mathsf{MATIME}^{\mathcal O}[2^n]\subseteq\mathrm{i.o.}\,\mathsf{SIZE}^{\mathcal O}[O(n)]$ for some oracle $\mathcal O$. We enumerate all the $2^n$-time randomized machine $M_1,M_2,\ldots$, where each $M_i$ takes an input $x$ of length $n$ and a certificate $y\in\{0,1\}^{2^n}$ and runs in time $2^n$. Note that any $\mathsf{MA}$-type algorithm that runs in time $2^n$ can be simulated by $M_i$ for some $i$.
Let $c$ be some universal constant ($c:=5$ suffices). Let $R_n:=\{0,1\}^{cn}$. The input to $\mathcal O$ is of the form $(r,i,x)$, where $r\in R_n$, $i\in[n]$, and $x\in\{0,1\}^n$ for some $n\in\mathbb N$. For any $r\in\{0,1\}^{cn}$, let $S_r$ be the set of $(r,i,x)$ such that $i\in[n]$ and $x\in\{0,1\}^n$.
We will construct an oracle $\mathcal O$ and a family of ``advice strings'' $r_n^*\in R_n$ for infinitely many $n$ such that
\begin{enumerate}
\item if $\Pr_{M_i}[M_i^{\mathcal O}(x,y)=1]\geq\frac34$ for some $y$, then $(r_n^*,i,x)\in\mathcal O$, and
\item if $\Pr_{M_i}[M_i^{\mathcal O}(x,y)=1]\leq\frac14$ for any $y$, then $(r_n^*,i,x)\notin\mathcal O$.
\end{enumerate}
Assuming this, it is easy to construct an $\mathcal O$-oracle linear-size circuit that simulates $M_i$ as follows. The circuit takes $r_n^*$ as hard-wired input and accepts an input $x\in\{0,1\}^n$ if and only if $(r_n,i,x)\in\mathcal O$.
Here is the construction of $\mathcal O$. We start with $\mathcal O:=\emptyset$. Some pairs $(i,x)$ will be marked ``forced'', meaning that $M_i$ accepts on input $x$. Some strings $r$ will be marked ``frozen'', meaning that the behavior of $\mathcal O$ on inputs in $S_r$ will not be changed after $r$ is frozen. Initially, no pair is forced and no advice string is frozen. We start with Stage $n:=1$. In Stage $n$, we construct an oracle as follows.
\paragraph{Stage $n$.} Consider the following condition, which we call $(*)_{\theta}$ for a threshold $\theta\in(0,1)$.
\begin{quote}
There exist an unfrozen string $r\in R_n$, an unforced pair $(i,x)\in[n]\times\{0,1\}^n$, an oracle $B\subseteq S_r$, and a certificate $y\in\{0,1\}^{2^n}$ such that
\[\Pr_{M_i}[M_i^{\mathcal O\cup B}(x,y)=1]\geq\theta,\]
where the probability is taken over the coin flip of the randomized machine $M_i$.
\end{quote}
We need to argue that the final oracle $\mathcal O$ can be computed in $\mathsf{EXPH}$. It may not be possible to check $(*)_{3/4}$ exactly in $\mathsf{EXPH}$, but using approximate counting in $\mathsf{PH}$ (Theorem23), we can check whether a promise variant of $(*)_{3/4}$ is satisfied or not in $\mathsf{EXPH}$. Specifically, by Theorem23, there exists an algorithm in $\mathsf{PH}$ that, given as input $n$ and $\mathcal O$ (which can be encoded as a binary string of length $2^{O(n)}$), accepts if $(*)_{3/4}$ holds, and rejects if $(*)_{1/2}$ does not hold. By the standard search-to-decision reduction, we obtain a $\mathsf{PH}$-oracle polynomial-time algorithm $S$ that outputs $r\in R_n$, $(i,x)\in[n]\times\{0,1\}^n$, $B\subseteq S_r$, and $y\in\{0,1\}^{2^n}$ that satisfy $(*)_{1/2}$ if $(*)_{3/4}$ holds, and outputs $\perp$ if $(*)_{1/2}$ does not hold. (Note that $S$ runs in time $2^{O(n)}$ on inputs of length $2^{O(n)}$ with a $\mathsf{PH}$ oracle.)
While the search algorithm $S$ outputs a certificate $(r,i,x,B)$ for $(*)_{1/2}$ on input $(n,\mathcal O)$ (instead of $\perp$), we do the following. Update $\mathcal O:=\mathcal O\cup B$. Let $r$ be frozen and let $(i,x)$ be forced. Let $R_{\theta}$ be the set of $r'\in R_n$ such that $(r',i',x')$ is queried during the computation of $M_i^{\mathcal O}(x,y)$ for some $(i',x')$ with probability at least $\theta$ over the internal randomness of $M_i$. By Theorem24, some set $A$ such that $R_{2^{-2n}}\subseteq A\subseteq\{0,1\}^*\setminus R_{2^{-2n-1}}$ can be enumerated in $\mathsf{EXPH}$. Note that $|A|\leq2^{2n+1}\cdot2^n$ because $M_i$ can make at most $2^n$ queries on each computation path. Let $r'$ be frozen for every $r'\in A$. If $S$ outputs $\perp$ (and thus $(*)_{3/4}$ does not hold), then let $r_n^*\in R_n$ be the first string that is not frozen. (We will later claim that such a string $r_n^*$ exists.) Then, we add to $\mathcal O$ all the tuples $(r_n^*,i,x)$ such that $(i,x)\in[n]\times\{0,1\}^n$ is forced. This completes the description of Stage $n$. Then we move on to the next stage $n':=2^n+1$, so that the construction in each stage is independent.
It is evident from the construction that $\mathcal O\in\mathsf{EXPH}$.
Fix any $n$ and consider Stage $n$. We claim that there exists a string $r_n^*\in R_n$ that is not frozen at the end of Stage $n$. We say that $(*)$ is satisfied if $S$ does not output $\perp$. Observe that the number of times that $(*)$ is satisfied is at most $n2^n$. The reason is that there are at most $n2^n$ pairs $(i,x)\in[n]\times\{0,1\}^n$, and each time $(*)$ is satisfied, $(i,x)$ becomes forced. Each time $(*)$ is satisfied, at most $1+2^{2n+1}\cdot2^n\leq2^{3n+2}$ strings $r\in R_n$ can be frozen. Thus, the number of frozen strings in $R_n$ in Stage $n$ is at most $n2^n\cdot2^{3n+2}$. Since there are at most $n$ stages before Stage $n$, in total, there are at most $n\cdot n2^n\cdot2^{3n+2}<2^{5n}\leq|R_n|$ frozen strings in $R_n$. (Here, we used that $c:=5$.) Thus, there exists some unfrozen string $r_n^*\in R_n$, which shows that the construction of $\mathcal O$ is well defined.
It remains to show that for every pair $(i,x)\in[n]\times\{0,1\}^n$,
\begin{enumerate}
\item if $\Pr_{M_i}[M_i^{\mathcal O}(x,y)=1]\geq\frac34$ for some $y$, then $(i,x)$ is forced (and thus $(r_n,i,x)\in\mathcal O$), and
\item if $\Pr_{M_i}[M_i^{\mathcal O}(x,y)=1]\leq\frac14$ for any $y$, then $(i,x)$ is not forced (and thus $(r_n,i,x)\notin\mathcal O$).
\end{enumerate}
Fix any $n$ and consider any unforced pair $(i,x)\in[n]\times\{0,1\}^n$. We claim that for every $y\in\{0,1\}^{2^n}$, it holds that
\[\Pr_{M_i}[M_i^{\mathcal O}(x,y)=1]<\tfrac34.\]
Fix any $y\in\{0,1\}^{2^n}$. Let $\mathcal O_n$ be the oracle $\mathcal O$ right after the while loop of Stage $n$. Since $(*)_{3/4}$ is not satisfied at that point, for every $r\in R_n$ and for every oracle $B\subseteq S_r$,
\[\Pr_{M_i}[M_i^{\mathcal O_n\cup B}(x,y)=1]<\tfrac34.\]
Let $B\subseteq S_{r_n^*}$ be the set of strings added to $\mathcal O$ at the end of Stage $n$. Then, the final oracle $\mathcal O$ coincides with $\mathcal O_n\cup B$ on any inputs of length at most $2^n$; thus, we obtain
\[\Pr_{M_i}[M_i^{\mathcal O}(x,y)=1]=\Pr_{M_i}[M_i^{\mathcal O_n\cup B}(x,y)=1]<\tfrac34.\]
Next, fix any $n$ and consider any forced pair $(i,x)\in[n]\times\{0,1\}^n$. Let $\mathcal O_{i,x}$ be the oracle right after $(i,x)$ is forced and $y$ be the certificate that satisfies $(*)$. Then, we have
\[\Pr[M_i^{\mathcal O_{i,x}}(x,y)=1]\geq\tfrac12.\]
We claim that this probability does not decrease by $\frac14$ even if we replace $\mathcal O_{i,x}$ with the final oracle $\mathcal O$. After $(i,x)$ is forced, $(*)$ can be satisfied at most $n2^n$ times, and each time $(*)$ is satisfied, at most one $S_r$ whose subset is added to $\mathcal O$. Thus, there are at most $n2^n$ unfrozen strings $r\in R_n$ such that some strings in $S_r$ are added to $\mathcal O$. Note that any string $r$ that is queried by $M_i$ with probability $2^{-2n}$ is frozen right after $(i,x)$ is forced. Thus, each unfrozen string can decrease the probability that $M_i^{\mathcal O}(x,y)$ accepts by $2^{-2n}$. In total, the unfrozen strings can affect the probability of acceptance by $n2^n\cdot2^{-2n}<\frac14$. Thus, we obtain
\[\Pr[M_i^{\mathcal O}(x,y)=1]>\tfrac12-\tfrac14\geq\tfrac14.\]
\end{proof}
